%==============================================================================
% HIGHLIGHTS — arquivo SEPARADO exigido pela GIQ (nome com "highlights").
% Regras: 3 a 5 itens; cada um com NO MÁXIMO 85 caracteres, contando espaços;
% devem captar os resultados novos e os métodos novos. Sem referências.
% A prosa é do autor. Abaixo, o que cada item deve dizer (um resultado por item),
% com as macros de numbers.tex que o sustentam. Compilar: pdflatex highlights.tex
%==============================================================================
\documentclass[12pt,a4paper]{article}
\usepackage[T1]{fontenc}
\usepackage[utf8]{inputenc}
\usepackage[margin=1in]{geometry}
\pagestyle{empty}
\begin{document}
\section*{Highlights}
\begin{itemize}
  \item [RESULT TO BE GENERATED] % 1. Achado central: a janela única muda o que o Estado registra (CobSantosPre -> CobSantosPos)
  \item [RESULT TO BE GENERATED] % 2. Tempos agregados: nenhum efeito robusto (TdoisPwcbFull, TquatroPwcbFull)
  \item [RESULT TO BE GENERATED] % 3. Efeito localizado: cabotagem recorrente, T4 (CabRecTquatroPct, CabRecTquatroPwcb)
  \item [RESULT TO BE GENERATED] % 4. Método: 23 portos datados por atos oficiais; DiD escalonado + bootstrap
  \item [RESULT TO BE GENERATED] % 5. Artefato: observatório reprodutível (pipeline + app) — opcional
\end{itemize}
\end{document}
